\section{Actuarial impact}
\label{sec:actuarial}

\noindent\textit{Skeleton; all tables are placeholders.}

\subsection{H5: propagation into life expectancy and annuity factors}

\begin{table}[htbp]\centering
\caption{Coverage of nominal 95\% intervals, interval width and error for
$e_0$, $e_{65}$ and $\annuity$ at 2\%, by family and mechanism, stable versus
shift regime.}
\label{tab:h5-actuarial}
\inputtable{tab-h5-actuarial}
\end{table}

The table will show whether an interval failure on the rate surface
survives aggregation into the quantities that are actually valued---or
whether the life-table integration averages it away. Both outcomes are
informative: the first says the reserve is wrong, the second says the
audit's headline is a statement about rates rather than about liabilities.
\todo{result pending}

\subsection{Reading the failure in balance-sheet terms}

For an annuity-due of one per annum at age 65, the interval on $\annuity$
translates directly into an interval on the best-estimate liability of a
unit annuitant. A nominal 95\% interval that covers the realised
$\annuity$ in materially fewer than 95\% of population--years is, in
Solvency~II language, an internal model whose longevity distribution is
too narrow; the size of the shortfall in the 97.5\% quantile of $\annuity$
is the amount by which the implied capital at that level is understated.
We report this quantity as a proportion of the central $\annuity$ so that
it is portable across populations and interest-rate assumptions.
\todo{decide whether to present a worked example on a hypothetical
portfolio; if so, no figures until the H5 table exists}

\subsection{The longevity tail}

The 99.5\% upper quantile of the predictive distribution of $\annuity$ is
the one-year longevity tail an internal model would carry. We report it
alongside the realised $\annuity$ for each cell; the fraction of cells in
which the realised value exceeds the model's own 99.5\% quantile is the
empirical frequency of a nominally once-in-200-year event, and is the
quantity the pandemic years let us observe for the first time.
\todo{descriptive; not a registered hypothesis, label as such}

\subsection{Twin crises in economic terms}

The placebo regime supplies the same functionals for 1914--1922. A
side-by-side of the $e_{65}$ and $\annuity$ interval shortfalls across the
two events asks whether the economic size of a break-induced interval
failure is stable across a century, or whether the pandemic years are
distinctive. \todo{result pending; gated on the placebo data-quality check}
